\documentclass[11pt,a4paper]{article}
\usepackage{amsmath,amssymb,booktabs,geometry}
\geometry{margin=2.4cm}
\title{An effect ledger for the continuous-time prompt\\on Game of 24}
\author{}
\date{2026}
\begin{document}
\maketitle

\section{What counts as an effect}
Each prompt quantity is compared with its own pre-drift mean.
The width of the gate is not part of the claim.
A quantity has an effect when its post-drift mean is higher than its pre-drift mean.
That rise is called an increment.
The business outcome \(Y\) is kept on its own line.
If the post-drift mean of \(Y\) is lower than its pre-drift mean, the outcome is recorded as a shortfall against the reference rate.

\(Y\) is the business label of a finished episode.
On Game of 24 it is \(1\) when search returns a state equal to \(24\), and \(0\) otherwise.
The same slot can hold a repayment, a completed order, or a closed ticket.
The residual, the online anomaly score, and the adjustment are written only after that label returns.
They do not exist on the query beat.

The clever covariate on a committed step is
\[
H=\frac{Y-e}{e(1-e)},
\]
where \(e\) is that step's stable score, clipped to \((0.001,0.999)\).
The online anomaly score is \(|H|\).
The sign is read separately: \(H>0.5\) is aggressive, \(H<-0.5\) is conservative, and the rest is keep.
With a binary \(Y\), a success gives \(H=1/e\) and a failure gives \(H=-1/(1-e)\).

\section{Game of 24}
Forty episodes, drift at episode 16, beam 4, depth 3.
The forest, the stack weight, and the feature split freeze after the first 12 episodes.
Reference \(Y\) is identically \(1\), so every column has permutation importance \(0\), the residual columns are empty, and the stack weight is \(0\).
The trail mean squared error is \(0.643\).
The out-of-bag mean squared error on the reference is \(0\).
On the 72 post-drift steps, the sign of \(H\) matches \(Y\) on every step.
The mean anomaly score is \(2.82\) on successes and \(1.24\) on failures.

\begin{center}
\begin{tabular}{lrrrl}
\toprule
Channel & Pre-drift & Post-drift & Change & Reading \\
\midrule
Terminal \(Y\) & 1.00 & 0.25 & \(-0.75\) & shortfall \\
Absolute residual & 0.00 & 0.75 & \(+0.75\) & increment \\
Online anomaly \(|H|\) & 2.38 & 1.63 & \(-0.75\) & below reference \\
Action prior \(e\) & 0.57 & 0.26 & \(-0.32\) & below reference \\
Conservative share & 0.00 & 0.75 & \(+0.75\) & increment \\
Hop rate & 0.00 & 0.25 & \(+0.25\) & increment \\
Last-step absolute error & 0.00 & 0.75 & \(+0.75\) & increment \\
Children scored per episode & 114.9 & 115.2 & \(+0.3\) & cost \\
Prompt window & 8 & 8 & 0 & cost \\
Prompt columns & 6 & 6 & 0 & cost \\
\bottomrule
\end{tabular}
\end{center}

\section{Why these four rises are the effect}
The forest stays frozen at a prediction of \(1\).
After \(Y\) is written back, the absolute residual on a step is \(1-Y\).
Its post-drift mean is therefore exactly the drop in the success rate, \(0.75\).
The same \(0.75\) appears as the conservative share, because every failed step has \(H<0\), and as the last-step absolute error.
The hop rate rises by \(0.25\): one quarter of the post-drift episodes move the successive squared error by enough to clear the hop.
These four quantities were zero on the pre-drift side and positive after the drift.
That is the effect.
How far above zero they sit is not a separate claim.

\(|H|\) does not rise.
The prior on this task stays below one half, so a success has the larger absolute value: \(1/e > 1/(1-e)\).
When failures become the majority, the average \(|H|\) falls from \(2.38\) to \(1.63\).
The anomaly score still matches the sign of the outcome, but the rise-above-reference test does not open on \(|H|\) for Game of 24.
The same split appeared on DoorKey (\(|H|\) from \(2.49\) to \(2.12\)) and the opposite split on the two-hop questions, where a prior clipped at \(0.001\) pulled the success score up to a post-drift mean of \(65.1\).

\section{What the increment is, and what it is not}
The increment on Game of 24 is the joint rise of the residual, the conservative share, the hop rate, and the last-step absolute error.
It is recorded on beats whose \(Y\) is already known, and it is visible to the next episode's prompt.
It does not raise the business outcome above the pre-drift rate.
That rate is already \(1\), and the post-drift rate is \(0.25\).
Scored children stay essentially flat, \(114.9\) to \(115.2\), because a depth-3 beam scores almost the same frontier whether or not the episode returns \(24\).
The window length and the column count do not change.

So the business line and the prompt line answer different questions.
\(Y\) says the solved-puzzle rate fell by \(0.75\) and did not clear the reference.
The prompt says four of its own quantities cleared theirs, by \(0.75\), \(0.75\), \(0.25\), and \(0.75\).
Both statements use the same episodes.

\end{document}
